\chapter{Duality and Reformulations}

\begin{abstract}
This chapter develops the Lagrange dual function, weak and strong duality,
dual decomposition, Wasserstein duality, epigraph and reparameterization
reformulations, and the standard convex classes LP, QP, QCQP, SOCP, and SDP.
\end{abstract}

\section{Duality and reformulations}
\label{sec:ch8-transformations}
The previous chapter asked a local question: given a constrained problem, what
must happen at a local minimizer? This chapter switches perspective. Instead
of extracting first-order information from the original problem, we reformulate
it: as a Lagrange dual problem, a split formulation, an epigraph form, or a
reparameterized unconstrained model. The goal is the same in every case:
preserve the essential optimization content, but move it into a representation
that is easier to analyze or solve.

\subsection{The dual function and the dual problem}
\label{subsec:ch8-dual}
Return to the primal problem
\begin{equation}
  \label{eq:ch8-primal-dual-primal}
  \begin{aligned}
    \min_{\mathbf{x}} \quad & f(\mathbf{x}) \\
    \text{subject to}\quad & g_i(\mathbf{x}) \le 0, \qquad i=1,\dots,m, \\
    & h_j(\mathbf{x}) = 0, \qquad j=1,\dots,p.
  \end{aligned}
\end{equation}
The associated Lagrangian is
\[
  L(\mathbf{x},\boldsymbol{\lambda},\boldsymbol{\mu})
  =
  f(\mathbf{x})
  + \sum_{i=1}^m \lambda_i g_i(\mathbf{x})
  + \sum_{j=1}^p \mu_j h_j(\mathbf{x}).
\]
The new object is the dual function
\begin{equation}
  \label{eq:ch8-dual-function}
  q(\boldsymbol{\lambda},\boldsymbol{\mu})
  :=
  \inf_{\mathbf{x}} L(\mathbf{x},\boldsymbol{\lambda},\boldsymbol{\mu}).
\end{equation}
For fixed multipliers, the Lagrangian is an objective in $\mathbf{x}$, and the
dual function records its best lower value.

\begin{claim}[Basic properties of the dual function]
\label{clm:ch8-dual-properties}
The function $q$ has the following properties.
\begin{enumerate}
  \item $q(\boldsymbol{\lambda},\boldsymbol{\mu})$ is concave in
  $(\boldsymbol{\lambda},\boldsymbol{\mu})$.
  \item For $ \forall \boldsymbol{\lambda}\ge \mathbf{0}, 
    \forall\boldsymbol{\mu}$
  \[
    q(\boldsymbol{\lambda},\boldsymbol{\mu}) \le f_{\mathrm{opt}},
  \]
  where $f_{\mathrm{opt}}$ is the optimal primal value.
\end{enumerate}
\end{claim}
\begin{proof}
For every fixed $\mathbf{x}$, the function
\[
  (\boldsymbol{\lambda},\boldsymbol{\mu})
  \mapsto
  L(\mathbf{x},\boldsymbol{\lambda},\boldsymbol{\mu})
\]
is affine, hence both convex and concave. Therefore
\[
  q(\boldsymbol{\lambda},\boldsymbol{\mu})
  =
  \inf_{\mathbf{x}} L(\mathbf{x},\boldsymbol{\lambda},\boldsymbol{\mu})
  =
  -\sup_{\mathbf{x}}
  \bigl(-L(\mathbf{x},\boldsymbol{\lambda},\boldsymbol{\mu})\bigr)
\]
is concave as the negative of a supremum of convex functions.

Now take any feasible $\widehat{\mathbf{x}}\in F$. Since
$g_i(\widehat{\mathbf{x}})\le 0$, $h_j(\widehat{\mathbf{x}})=0$, and
$\lambda_i\ge 0$, we obtain
\begin{align*}
  q(\boldsymbol{\lambda},\boldsymbol{\mu})
  &=
  \inf_{\mathbf{x}} L(\mathbf{x},\boldsymbol{\lambda},\boldsymbol{\mu}) \\
  &\le
  L(\widehat{\mathbf{x}},\boldsymbol{\lambda},\boldsymbol{\mu}) \\
  &=
  f(\widehat{\mathbf{x}})
  + \sum_{i=1}^m \lambda_i g_i(\widehat{\mathbf{x}})
  + \sum_{j=1}^p \mu_j h_j(\widehat{\mathbf{x}}) \\
  &\le f(\widehat{\mathbf{x}}).
\end{align*}
Taking the infimum over all feasible $\widehat{\mathbf{x}}$ gives
\[
  q(\boldsymbol{\lambda},\boldsymbol{\mu}) \le f_{\mathrm{opt}}.
\]
\end{proof}

Because $q$ is concave and the constraint
$\boldsymbol{\lambda}\ge \mathbf{0}$ is convex, the dual maximization has the
standard convex form: equivalently, one minimizes the convex function $-q$ over
a convex feasible set. This gives the dual problem
\begin{equation}
  \label{eq:ch8-dual-problem}
  \begin{aligned}
    \max_{\boldsymbol{\lambda},\boldsymbol{\mu}} \quad
    & q(\boldsymbol{\lambda},\boldsymbol{\mu}) \\
    \text{subject to}\quad
    & \boldsymbol{\lambda}\ge \mathbf{0}.
  \end{aligned}
\end{equation}

\begin{theorem}[Weak duality]
\label{thm:ch8-weak-duality}
For every primal feasible $\mathbf{x}$ and every dual feasible
$(\boldsymbol{\lambda},\boldsymbol{\mu})$,
\[
  q(\boldsymbol{\lambda},\boldsymbol{\mu}) \le f(\mathbf{x}).
\]
In particular, if $q_{\mathrm{opt}}$ denotes the optimal dual value, then
\[
  q_{\mathrm{opt}} \le f_{\mathrm{opt}}.
\]
\end{theorem}

The quantity
\[
  f_{\mathrm{opt}} - q_{\mathrm{opt}} \ge 0
\]
is the \emph{duality gap}. The main favorable case is strong duality.

\begin{theorem}[Strong duality in the convex case]
\label{thm:ch8-strong-duality}
Assume that, in the primal problem \eqref{eq:ch8-primal-dual-primal}, the
objective $f$ and the inequality functions $g_i$ are convex, the equality
functions $h_j$ are affine, and Slater's condition from \Cref{def:ch7-cq}
holds. Then
\[
  q_{\mathrm{opt}} = f_{\mathrm{opt}}.
\]
If, in addition, dual optimizers
$(\boldsymbol{\lambda}^\star,\boldsymbol{\mu}^\star)$ exist, then every primal
optimizer $\mathbf{x}^\star$ satisfies
\[
  \mathbf{x}^\star \in
  \operatorname*{Argmin}_{\mathbf{x}}
  L(\mathbf{x},\boldsymbol{\lambda}^\star,\boldsymbol{\mu}^\star).
\]
\end{theorem}

This explains why the dual problem is useful. First, the dual objective
$q$ is always concave, so the dual can be viewed as a convex optimization
problem by minimizing $-q$ over the convex set
$\{(\boldsymbol{\lambda},\boldsymbol{\mu}):\boldsymbol{\lambda}\ge\mathbf{0}\}$.
This can be more pleasant to solve than the primal problem, especially when
the primal is not convex or has awkward constraints. Second, dual formulations
often expose structure that is hidden in the primal. A central example is the
support vector machine: after deriving its dual, the data enter only through
inner products, which is exactly what makes the kernel trick possible.

\begin{example}[Dual of a linear program with equalities and inequalities]
\label{ex:ch8-dual-lp}
Consider
\[
  \min_{\mathbf{x}} c^\top \mathbf{x}
  \qquad
  \text{subject to }
  A\mathbf{x}\le \mathbf{b},
  \qquad
  C\mathbf{x}=\mathbf{d}.
\]
Its Lagrangian is
\[
  L(\mathbf{x},\boldsymbol{\lambda},\boldsymbol{\mu})
  =
  c^\top \mathbf{x}
  + \boldsymbol{\lambda}^\top(A\mathbf{x}-\mathbf{b})
  + \boldsymbol{\mu}^\top(C\mathbf{x}-\mathbf{d}),
\]
or, after grouping the $\mathbf{x}$-terms,
\[
  L(\mathbf{x},\boldsymbol{\lambda},\boldsymbol{\mu})
  =
  \mathbf{x}^\top
  (c + A^\top \boldsymbol{\lambda} + C^\top \boldsymbol{\mu})
  - \boldsymbol{\lambda}^\top \mathbf{b}
  - \boldsymbol{\mu}^\top \mathbf{d}.
\]
Therefore
\[
  q(\boldsymbol{\lambda},\boldsymbol{\mu})
  =
  \inf_{\mathbf{x}} L(\mathbf{x},\boldsymbol{\lambda},\boldsymbol{\mu})
  =
  \begin{cases}
    -\boldsymbol{\lambda}^\top \mathbf{b}
    - \boldsymbol{\mu}^\top \mathbf{d},
    &
    c + A^\top \boldsymbol{\lambda} + C^\top \boldsymbol{\mu}
    = \mathbf{0},
    \\
    -\infty, & \text{otherwise}.
  \end{cases}
\]
The reason is simple: if the coefficient of $\mathbf{x}$ is nonzero, then the
affine function of $\mathbf{x}$ can be driven to $-\infty$.

Hence the dual problem is
\[
  \max_{\boldsymbol{\lambda},\boldsymbol{\mu}}
  \bigl(
    -\boldsymbol{\lambda}^\top \mathbf{b}
    - \boldsymbol{\mu}^\top \mathbf{d}
  \bigr)
  \qquad
  \text{subject to }
  \boldsymbol{\lambda}\ge \mathbf{0},
  \qquad
  c + A^\top \boldsymbol{\lambda} + C^\top \boldsymbol{\mu}=\mathbf{0}.
\]
This example also shows a limitation of \Cref{thm:ch8-strong-duality}. When
$c + A^\top \boldsymbol{\lambda} + C^\top \boldsymbol{\mu}= \mathbf{0}$, the
Lagrangian is constant in $\mathbf{x}$, so
$\operatorname*{Argmin}_{\mathbf{x}}
L(\mathbf{x},\boldsymbol{\lambda},\boldsymbol{\mu})
=\mathbb{R}^n$. Thus the theorem does not identify the true primal optimizer;
it only says that any true primal optimizer lies among these Lagrangian
minimizers.
\end{example}

\subsection{Dual decomposition}
\label{subsec:ch8-dual-decomposition}
A standard splitting idea starts from an objective that is a
sum of two simple pieces, but optimizing them together is inconvenient:
\[
  \min_{\mathbf{x}} f(\mathbf{x}) + g(\mathbf{x}).
\]
Introduce two copies of the variable and tie them together by an equality:
\begin{equation}
  \label{eq:ch8-dual-decomposition-primal}
  \min_{\mathbf{x},\mathbf{y}} f(\mathbf{x}) + g(\mathbf{y})
  \qquad
  \text{subject to }
  \mathbf{x}=\mathbf{y}.
\end{equation}
Its Lagrangian is
\[
  L(\mathbf{x},\mathbf{y},\boldsymbol{\mu})
  =
  f(\mathbf{x}) + g(\mathbf{y}) + \boldsymbol{\mu}^\top(\mathbf{x}-\mathbf{y}),
\]
so the dual function separates:
\begin{align}
  q(\boldsymbol{\mu})
  &=
  \inf_{\mathbf{x},\mathbf{y}}
  L(\mathbf{x},\mathbf{y},\boldsymbol{\mu}) \notag \\
  &=
  \inf_{\mathbf{x}}
  \bigl(f(\mathbf{x})+\boldsymbol{\mu}^\top \mathbf{x}\bigr)
  +
  \inf_{\mathbf{y}}
  \bigl(g(\mathbf{y})-\boldsymbol{\mu}^\top \mathbf{y}\bigr).
  \label{eq:ch8-dual-decomposition-q}
\end{align}
This is the basic advantage of the reformulation: the primal coupling
$\mathbf{x}=\mathbf{y}$ becomes an additive separation in the dual.

If the minimizers
\[
  \mathbf{x}_{\mathrm{opt}}(\boldsymbol{\mu})
  \in
  \arg\min_{\mathbf{x}}
  \bigl(f(\mathbf{x})+\boldsymbol{\mu}^\top \mathbf{x}\bigr),
  \qquad
  \mathbf{y}_{\mathrm{opt}}(\boldsymbol{\mu})
  \in
  \arg\min_{\mathbf{y}}
  \bigl(g(\mathbf{y})-\boldsymbol{\mu}^\top \mathbf{y}\bigr)
\]
are well defined, then the gradient of the dual function is
\[
  \nabla_{\boldsymbol{\mu}} q(\boldsymbol{\mu})
  =
  \mathbf{x}_{\mathrm{opt}}(\boldsymbol{\mu})
  -
  \mathbf{y}_{\mathrm{opt}}(\boldsymbol{\mu}).
\]
So gradient methods can be applied on the dual side. At each multiplier
$\boldsymbol{\mu}$, first solve the two separated primal subproblems to find
$\mathbf{x}_{\mathrm{opt}}(\boldsymbol{\mu})$ and
$\mathbf{y}_{\mathrm{opt}}(\boldsymbol{\mu})$. Then use the gradient above in
a gradient ascent method for $q$ with respect to $\boldsymbol{\mu}$.

\begin{example}[$\ell_1$-regularized least squares]
\label{ex:ch8-dual-lasso}
Consider the problem
\[
  \min_{\mathbf{x}}
  \frac12 \|A\mathbf{x}-\mathbf{b}\|_2^2 + \lambda \|\mathbf{x}\|_1.
\]
The nonsmooth term and the quadratic term become easier to separate after
introducing
\[
  \mathbf{z}=A\mathbf{x}-\mathbf{b}.
\]
This gives the equivalent constrained problem
\[
  \min_{\mathbf{x},\mathbf{z}}
  \frac12 \|\mathbf{z}\|_2^2 + \lambda \|\mathbf{x}\|_1
  \qquad
  \text{subject to }
  A\mathbf{x}-\mathbf{b}-\mathbf{z}=\mathbf{0}.
\]
Its Lagrangian is
\[
  L(\mathbf{x},\mathbf{z},\boldsymbol{\mu})
  =
  \frac12 \|\mathbf{z}\|_2^2
  + \lambda \|\mathbf{x}\|_1
  + \boldsymbol{\mu}^\top(A\mathbf{x}-\mathbf{b}-\mathbf{z}).
\]
Minimizing first with respect to $\mathbf{z}$ gives
\[
  \nabla_{\mathbf{z}} L = \mathbf{z}-\boldsymbol{\mu}=\mathbf{0},
  \qquad
  \mathbf{z}=\boldsymbol{\mu},
\]
and hence
\[
  \inf_{\mathbf{z}}
  \left(
    \frac12 \|\mathbf{z}\|_2^2
    - \boldsymbol{\mu}^\top \mathbf{z}
  \right)
  =
  -\frac12 \|\boldsymbol{\mu}\|_2^2.
\]
For the $\mathbf{x}$-part,
\[
  \lambda \|\mathbf{x}\|_1 + \boldsymbol{\mu}^\top A\mathbf{x}
  =
  \sum_i
  \Bigl(
    \lambda |x_i| + (A^\top \boldsymbol{\mu})_i x_i
  \Bigr).
\]
Each scalar term has finite infimum if and only if
\[
  -\lambda \le (A^\top \boldsymbol{\mu})_i \le \lambda,
  \qquad i=1,\dots,n,
\]
in which case the infimum equals $0$; otherwise the infimum is $-\infty$.
Therefore the dual problem is
\[
  \max_{\boldsymbol{\mu}}
  \left(
    -\frac12 \|\boldsymbol{\mu}\|_2^2
    - \boldsymbol{\mu}^\top \mathbf{b}
  \right)
  \qquad
  \text{subject to }
  \|A^\top \boldsymbol{\mu}\|_\infty \le \lambda.
\]
This is now a standard quadratic maximization problem
with simple convex constraints.
After maximizing with respect to $\boldsymbol{\mu}$, the optimal multiplier
$\boldsymbol{\mu}^\star$ gives
$\mathbf{z}^\star=\boldsymbol{\mu}^\star$ from the stationarity condition in
$\mathbf{z}$. The constraint
$A\mathbf{x}^\star-\mathbf{b}=\mathbf{z}^\star$ then gives
\[
  A\mathbf{x}^\star=\mathbf{b}+\boldsymbol{\mu}^\star.
\]
When this linear system determines $\mathbf{x}^\star$ uniquely, it recovers the
primal optimizer.
\end{example}

\subsection{Wasserstein distance as a duality example}
\label{subsec:ch8-wasserstein}
A particularly useful duality example comes from optimal transport and its
machine-learning applications. The physical picture is the following. Think of
$p$ as a pile of unit mass spread over space, and think of $q$ as the desired
final arrangement of the same mass. To turn $p$ into $q$, we choose how much
mass to move from each point $x$ to each point $y$. Moving a large amount of
mass is expensive, and moving it far is also expensive; the cost of sending
mass from $x$ to $y$ is therefore
\[
  \text{mass moved from } x \text{ to } y
  \times
  \|x-y\|.
\]
The Wasserstein distance is the minimum total work needed to rearrange the mass
of $p$ into the mass of $q$. This is why it is also called the
\emph{earth mover's distance}.

This point of view is especially useful when the supports of the two
distributions do not overlap. Pointwise divergences such as KL divergence
compare densities at the same location and can become infinite when
$q(x)=0$ on a region where $p(x)>0$. 

For discrete distributions $p$ and $q$, a transport plan is a nonnegative
function $d(x,y)$ satisfying the marginal constraints
\[
  d(x,y)\ge 0,
  \qquad
  \sum_y d(x,y)=p(x),
  \qquad
  \sum_x d(x,y)=q(y).
\]
The corresponding optimal transport problem is
\[
  \min_d \sum_{x,y} d(x,y)\,\|x-y\|
  \qquad
  \text{subject to the constraints above.}
\]

In continuous form, the Wasserstein distance is
\[
  W(p,q)
  =
  \inf_{d\in \Pi(p,q)}
  \mathbb{E}_{(x,y)\sim d}\|x-y\|,
\]
where $\Pi(p,q)$ is the set of all couplings with marginals $p$ and $q$:
\[
  d(x,y)\ge 0,
  \qquad
  \int d(x,y)\,dy = p(x),
  \qquad
  \int d(x,y)\,dx = q(y).
\]

\begin{theorem}[Kantorovich--Rubinstein duality]
\label{thm:ch8-kr}
The Wasserstein distance satisfies
\[
  W(p,q)
  =
  \sup_{h\in 1\text{-Lipschitz}}
  \left(
    \mathbb{E}_{x\sim p} h(x)
    -
    \mathbb{E}_{y\sim q} h(y)
  \right).
\]
\end{theorem}
\begin{proof}
Start from the constrained primal problem
\[
  \inf_{d}
  \iint d(x,y)\|x-y\|\,dx\,dy
  \quad
  \text{subject to }
  \begin{cases}
    d(x,y)\ge 0, \\
    \int d(x,y)\,dy = p(x), \\
    \int d(x,y)\,dx = q(y).
  \end{cases}
\]
Introduce dual functions $\mu_1(x)$ and $\mu_2(y)$ for the marginal
constraints and a nonnegative multiplier $\lambda(x,y)$ for $d(x,y)\ge 0$.
The Lagrangian is
\begin{align*}
  L[d,\mu_1,\mu_2,\lambda]
  &=
  \iint d(x,y)\|x-y\|\,dx\,dy \\
  &\quad +
  \int \mu_1(x)
  \left(
    p(x) - \int d(x,y)\,dy
  \right)
  dx \\
  &\quad +
  \int \mu_2(y)
  \left(
    q(y) - \int d(x,y)\,dx
  \right)
  dy \\
  &\quad -
  \iint \lambda(x,y)d(x,y)\,dx\,dy.
\end{align*}
After collecting the $d$-terms, this becomes
\begin{align*}
  L[d,\mu_1,\mu_2,\lambda]
  &=
  \iint d(x,y)
  \bigl(
    \|x-y\|-\mu_1(x)-\mu_2(y)-\lambda(x,y)
  \bigr)
  dx\,dy \\
  &\quad +
  \mathbb{E}_{x\sim p}\mu_1(x)
  +
  \mathbb{E}_{y\sim q}\mu_2(y).
\end{align*}
So the dual function is finite only if
\[
  \|x-y\|-\mu_1(x)-\mu_2(y)-\lambda(x,y)=0,
  \qquad
  \lambda(x,y)\ge 0,
\]
that is,
\[
  \mu_1(x)+\mu_2(y)\le \|x-y\|
  \qquad \forall x,y.
\]
Hence the dual problem is
\begin{equation}
  \label{eq:ch8-wasserstein-dual-raw}
  \sup_{\mu_1,\mu_2}
  \left(
    \mathbb{E}_{x\sim p}\mu_1(x)
    +
    \mathbb{E}_{y\sim q}\mu_2(y)
  \right)
  \quad
  \text{subject to }
  \mu_1(x)+\mu_2(y)\le \|x-y\|
  \ \forall x,y.
\end{equation}
Here we use strong duality for the optimal-transport linear program: the
Wasserstein distance $W(p,q)$ equals the value of
\eqref{eq:ch8-wasserstein-dual-raw}.\footnote{This strong-duality result for
optimal transport is not proved here.} 

Now let $h$ be any $1$-Lipschitz function. For every coupling $d\in \Pi(p,q)$,
\begin{align*}
  \mathbb{E}_{x\sim p} h(x) - \mathbb{E}_{y\sim q} h(y)
  &=
  \iint d(x,y)\bigl(h(x)-h(y)\bigr)\,dx\,dy \\
  &\le
  \iint d(x,y)\|x-y\|\,dx\,dy.
\end{align*}
Taking the infimum over all couplings gives
\[
  \mathbb{E}_{x\sim p} h(x) - \mathbb{E}_{y\sim q} h(y)
  \le W(p,q),
\]
and then taking the supremum over all $1$-Lipschitz $h$ yields
\begin{equation}
  \label{eq:ch8-kr-one-side}
  \sup_{h\in 1\text{-Lipschitz}}
  \left(
    \mathbb{E}_{x\sim p} h(x)
    -
    \mathbb{E}_{y\sim q} h(y)
  \right)
  \le W(p,q).
\end{equation}

For the reverse inequality, take any feasible pair $(\mu_1,\mu_2)$ in
\eqref{eq:ch8-wasserstein-dual-raw} and define
\[
  k(x)
  :=
  \inf_y \bigl(\|x-y\|-\mu_2(y)\bigr).
\]
Then \(\mu_1(x)\le k(x)\) \(\forall x\). Also, for any $x,z$,
\[
  k(x)
  \le
  \|x-z\|
  +
  \inf_y \bigl(\|z-y\|-\mu_2(y)\bigr)
  =
  \|x-z\| + k(z),
\]
so
\[
  |k(x)-k(z)|\le \|x-z\|,
\]
which means that $k$ is $1$-Lipschitz. Finally, by definition of $k$, if we
take $y=x$, then
\[
  k(x)\le \|x-x\|-\mu_2(x) = -\mu_2(x),
\]
or equivalently $\mu_2(x)\le -k(x)$. Therefore
\[
  \mathbb{E}_{x\sim p}\mu_1(x)
  +
  \mathbb{E}_{y\sim q}\mu_2(y)
  \le
  \mathbb{E}_{x\sim p}k(x)
  -
  \mathbb{E}_{y\sim q}k(y)
  \le
  \sup_{h\in 1\text{-Lipschitz}}
  \left(
    \mathbb{E}_{x\sim p} h(x)
    -
    \mathbb{E}_{y\sim q} h(y)
  \right),
\]
where the last inequality holds because $k$ is $1$-Lipschitz. Taking the
supremum over feasible $(\mu_1,\mu_2)$ and using strong duality for the raw
dual problem gives
\[
  W(p,q)
  =
  \sup_{(\mu_1,\mu_2)\ \text{feasible}}
  \left(
    \mathbb{E}_{x\sim p}\mu_1(x)
    +
    \mathbb{E}_{y\sim q}\mu_2(y)
  \right)
  \le
  \sup_{h\in 1\text{-Lipschitz}}
  \left(
    \mathbb{E}_{x\sim p} h(x)
    -
    \mathbb{E}_{y\sim q} h(y)
  \right).
\]
Together with \eqref{eq:ch8-kr-one-side}, this gives
\[
  W(p,q)
  \le
  \sup_{h\in 1\text{-Lipschitz}}
  \left(
    \mathbb{E}_{x\sim p} h(x)
    -
    \mathbb{E}_{y\sim q} h(y)
  \right)
  \le
  W(p,q),
\]
so equality holds.
\end{proof}

This duality formula is one of the reasons Wasserstein distance is useful in
machine learning: it replaces optimization over couplings by optimization over
functions with a Lipschitz constraint.

\begin{example}[Wasserstein GAN]
\label{ex:ch8-wgan}
Suppose we are given samples
\[
  X_1,\dots,X_N,
\]
with empirical data distribution
\[
  p_{\mathrm{data}}(x)
  =
  \frac1N \sum_{i=1}^N \delta(x-X_i).
\]
Let $z\sim p(z)$ be a latent variable and let the generator produce
\[
  x = g(z,\theta),
\]
which induces a model distribution $q_\theta$. By
\Cref{thm:ch8-kr}, minimizing the Wasserstein distance can be written as
\[
  \min_\theta
  \sup_{h\in 1\text{-Lipschitz}}
  \left(
    \mathbb{E}_{x\sim p_{\mathrm{data}}} h(x)
    -
    \mathbb{E}_{z\sim p(z)} h(g(z,\theta))
  \right).
\]
In practice we model the $1$-Lipschitz function by a neural network,
\[
  h(\cdot,w)=\text{neural network},
\]
with weights $w$. Unlike the discriminator in the original GAN formulation,
this network is called the \emph{critic}.

The remaining issue is the $1$-Lipschitz constraint. One approach is to
restrict the size of the network weights. For a linear layer,
\[
  \|W x_1-W x_2\|
  \le
  \|W\|\|x_1-x_2\|,
\]
so its Lipschitz constant is controlled by $\|W\|$. For a network
which is a composition of linear layers and $1$-Lipschitz activations,
\[
  h
  =
  W_L\circ \sigma_{L-1}\circ W_{L-1}\circ \cdots \circ \sigma_1\circ W_1,
  \qquad
  \operatorname{Lip}(h)
  \le
  \prod_{\ell=1}^L \|W_\ell\|.
\]
Thus controlling the sizes of the weight matrices helps control
$\operatorname{Lip}(h)$. 

Another approach is to add a gradient penalty term of the form
\[
  \lambda\bigl(\|\nabla h(x)\|_2 - 1\bigr)^2
\]
to the critic objective. By the multivariate mean value theorem, for
differentiable $h$ and some point $\xi$ on the segment between $x_1$ and
$x_2$,
\[
  |h(x_1)-h(x_2)|
  =
  |\nabla h(\xi)^\top(x_1-x_2)|
  \le
  \|\nabla h(\xi)\|\|x_1-x_2\|.
\]
Hence the local condition
\[
  \|\nabla h(x)\|\le 1
\]
implies the global $1$-Lipschitz condition. 
\end{example}

\subsection{Epigraph reformulations}
\label{subsec:ch8-epigraph}
Chapter~4 already introduced epigraphs as geometric objects. Here the same
idea appears as a modeling trick. The basic identity is
\[
  \min_{\mathbf{x}} f(\mathbf{x})
  \qquad \Longleftrightarrow \qquad
  \min_{\mathbf{x},t} t
  \quad
  \text{subject to }
  f(\mathbf{x})\le t.
\]
The objective has been linearized, at the cost of one extra constraint.

\begin{example}[The infinity norm as a mini-max problem]
\label{ex:ch8-epigraph-inf}
Consider
\[
  \min_{\mathbf{x}} \|A\mathbf{x}-\mathbf{b}\|_\infty
  =
  \min_{\mathbf{x}} \max_i |(A\mathbf{x}-\mathbf{b})_i|.
\]
Introducing an epigraph variable $t$ gives
\[
  \min_{\mathbf{x},t} t
  \qquad
  \text{subject to }
  |(A\mathbf{x}-\mathbf{b})_i|\le t
  \quad \forall i.
\]
Equivalently,
\[
  \min_{\mathbf{x},t} t
  \qquad
  \text{subject to }
  -t \le (A\mathbf{x}-\mathbf{b})_i \le t
  \quad \forall i.
\]
So a mini-max objective becomes a linear objective with linear inequalities.
\end{example}

\begin{example}[Turning an $\ell_1$ term into linear constraints]
\label{ex:ch8-epigraph-l1}
For
\[
  \min_{\mathbf{x}} f(\mathbf{x}) + \lambda \|\mathbf{x}\|_1,
\]
introduce one slack variable $t_j$ for each coordinate. Then the equivalent
problem is
\[
  \min_{\mathbf{x},\mathbf{t}\in \mathbb{R}^d}
  f(\mathbf{x}) + \lambda \sum_{j=1}^d t_j
  \qquad
  \text{subject to }
  -t_j \le x_j \le t_j,
  \quad j=1,\dots,d.
\]
The nonsmooth absolute values have disappeared from the objective and been
moved into linear constraints.
\end{example}

\subsection{Reparameterization reformulations}
\label{subsec:ch8-reparam}
Another useful idea is to enforce the constraint set by construction.
If $\Omega$ is the feasible set and we can write
\[
  \mathbf{x}=g(\mathbf{y})\in \Omega,
  \qquad
  \mathbf{y}\in \mathbb{R}^m,
\]
then
\[
  \min_{\mathbf{x}\in \Omega} f(\mathbf{x})
  \qquad \Longleftrightarrow \qquad
  \min_{\mathbf{y}} f(g(\mathbf{y})).
\]
This removes the explicit constraint, but it changes the geometry of the
problem and may create new numerical issues.

Typical examples are the following.
\begin{enumerate}
  \item For the scalar constraint $x\ge 0$, use
  \[
    x = g(y) = \log(1+e^y).
  \]
  The seemingly simpler map $x=y^2$ is often a
  poor choice numerically: it is symmetric, not monotone, and may create bad
  gradients.

  \item For the box constraint $a\le x\le b$, use
  \[
    x = g(y) = a + (b-a)\sigma(y),
  \]
  where $\sigma(y)=1/(1+e^{-y})$. 

  \item For positive semidefinite matrices,
  \[
    X \succeq 0
    \qquad \Longrightarrow \qquad
    X = Y Y^\top.
  \]

  \item For positive definite matrices, use a Cholesky-type factorization
  \[
    X = L L^\top,
  \]
  where $L$ is triangular and its diagonal is forced to be positive, for
  example by
  \[
    L_{ii} = \log(1+e^{\ell_{ii}}).
  \]

  \item For orthogonality constraints,
  \[
    X^\top X = I
    \qquad \Longrightarrow \qquad
    X = \exp(Y-Y^\top).
  \]
  This reaches the connected component with determinant $1$.

  \item For the Euclidean ball $\|\mathbf{x}\|_2\le a$, one can separate radius
  and direction, for instance by
  \[
    \mathbf{x}
    =
    a\,\delta(z)\,
    \frac{\mathbf{y}}{\|\mathbf{y}\|_2},
  \]
  where $\delta(z)\in [0,1]$. Another useful map is the
  projection-style parameterization
  \[
    g(\mathbf{y})
    =
    \operatorname{proj}_{\{\|\mathbf{u}\|_2\le a\}}(\mathbf{y})
    =
    \begin{cases}
      \mathbf{y}, & \|\mathbf{y}\|_2\le a, \\
      a\,\mathbf{y}/\|\mathbf{y}\|_2, & \text{otherwise}.
    \end{cases}
  \]
\end{enumerate}

\subsection{Standard classes of constrained convex optimization}
\label{subsec:ch8-standard-classes}
We close with the model classes that appear repeatedly in
applications. They all fit the constrained template of this chapter, but each
class has specialized algorithms and geometry.

\paragraph{Linear programming (LP).}
\[
  \min_{\mathbf{x}} c^\top \mathbf{x}
  \qquad
  \text{subject to }
  A\mathbf{x}\le \mathbf{b},
  \qquad
  E\mathbf{x}=\mathbf{f}.
\]

\paragraph{Quadratic programming (QP).}
\[
  \min_{\mathbf{x}}
  \frac12 \mathbf{x}^\top A \mathbf{x} + \mathbf{x}^\top \mathbf{b}
  \qquad
  \text{subject to }
  C\mathbf{x}\le \mathbf{d},
  \qquad
  E\mathbf{x}=\mathbf{f},
  \qquad
  A\in \mathbb{S}_+^n.
\]

\paragraph{Quadratically constrained quadratic programming (QCQP).}
\[
  \min_{\mathbf{x}}
  \frac12 \mathbf{x}^\top A \mathbf{x} + \mathbf{x}^\top \mathbf{b}
  \qquad
  \text{subject to }
  \frac12 \mathbf{x}^\top P_j \mathbf{x}
  + \mathbf{x}^\top \mathbf{q}_j + r_j \le 0,
  \quad j=1,\dots,m,
\]
\[
  E\mathbf{x}=\mathbf{f},
  \qquad
  A,P_j\in \mathbb{S}_+^n.
\]

\paragraph{Second-order cone programming (SOCP).}
\[
  \min_{\mathbf{x}} c^\top \mathbf{x}
  \qquad
  \text{subject to }
  \|A_i\mathbf{x}+\mathbf{b}_i\|_2 \le c_i^\top \mathbf{x} + d_i,
  \quad i=1,\dots,m,
  \qquad
  E\mathbf{x}=\mathbf{f}.
\]

\paragraph{Semidefinite programming (SDP).}
\[
  \min_{\mathbf{x}\in \mathbb{R}^n} c^\top \mathbf{x}
  \qquad
  \text{subject to }
  A_0 + x_1 A_1 + \cdots + x_n A_n \succeq 0,
  \qquad
  E\mathbf{x}=\mathbf{f},
\]
with $A_i\in \mathbb{S}^m$.

These classes satisfy the inclusion chain
\[
  \mathrm{LP} \subset \mathrm{QP} \subset \mathrm{QCQP}
  \subset \mathrm{SOCP} \subset \mathrm{SDP}.
\]
As Chapter~4 already suggested through epigraphs, cones, and positive
semidefinite matrices, each inclusion means that the later class can represent
the earlier one.

